\section*{Chapter \ref{ch:b2:atomic-orbitals} --- Atomic Orbitals: Shapes, Energies and Sizes}

\begin{solution}{exo:b2:atomic-orbitals:1}
3p: $n = 3$, $l = 1$; one \omterm{def:b2:atomic-orbitals:node}{radial node} ($n - l - 1$) and one \omterm{def:b2:atomic-orbitals:node}{angular node}. 4d: $n = 4$,
$l = 2$; one \omterm{def:b2:atomic-orbitals:node}{radial node} and two \omterm{def:b2:atomic-orbitals:node}{angular nodes}.
\end{solution}

\begin{solution}{exo:b2:atomic-orbitals:2}
$P_{2p} \propto r^4\mathrm e^{-r}$ ($r$ in $a_0$): $d\ln P/dr = 4/r - 1 = 0$ for $r = 4$:
$4a_0 = \qty{212}{pm}$.
\end{solution}

\begin{solution}{exo:b2:atomic-orbitals:3}
$1 - \mathrm e^{-2}(1 + 2 + 2) = 1 - 5\mathrm e^{-2} = 0.323$.
\end{solution}

\begin{solution}{exo:b2:atomic-orbitals:4}
Four lobes pointing between the $x$ and $y$ axes, in the $xy$ plane, of alternating
signs (positive where $xy > 0$); the nodal planes are $xz$ and $yz$. \omterm{def:b2:atomic-orbitals:node}{Radial nodes}:
$n - l - 1 = 0$.
\end{solution}

\begin{solution}{exo:b2:atomic-orbitals:5}
Oxygen, $(1\mathrm s)^2(2\mathrm s, 2\mathrm p)^6$: $\sigma = 5 \times 0.35 + 2 \times 0.85
= 3.45$, $Z^* = 4.55$; radius $4/4.55 = 0.88a_0$, \qty{47}{pm}.
\end{solution}

\begin{solution}{exo:b2:atomic-orbitals:6}
Ionisation removes the least bound electrons. The 4s electrons
(\qty{-14}{eV}) are far less bound than the 3d ones (\qty{-59}{eV}): both 4s
electrons leave first, then one 3d electron, giving $[\ce{Ar}]\,3\mathrm d^5$.
\end{solution}

\begin{solution}{exo:b2:atomic-orbitals:7}
Lithium, 2s: $Z^* = 3 - 2 \times 0.85 = 1.30$, $I = 13.6 \times 1.30^2/4 =
\qty{5.75}{eV}$ (measured 5.39, 7~\% too high). Sodium, 3s: $Z^* = 11 - 8 \times 0.85 -
2 \times 1.00 = 2.20$, $I = 13.6 \times 2.20^2/9 = \qty{7.3}{eV}$ (measured 5.14): the
rules, a rough model, overestimate the binding of the outer electron, more and
more down the group.
\end{solution}

\begin{solution}{exo:b2:atomic-orbitals:8}
Sodium: $Z^* = 2.20$, radius $9/2.20 = 4.1a_0$ (\qty{216}{pm}). Chlorine, 3p:
$\sigma = 6 \times 0.35 + 8 \times 0.85 + 2 \times 1.00 = 10.90$, $Z^* = 6.10$, radius
$9/6.10 = 1.5a_0$ (\qty{78}{pm}). Same shell, but the outer electron of chlorine
feels almost three times the charge, and is pulled in.
\end{solution}

\begin{solution}{exo:b2:atomic-orbitals:9}
Both have the \omterm{def:b2:atomic-orbitals:radial-angular}{angular part} $1/\sqrt{4\pi}$, so the integral reduces to
\[
  \int_0^\infty R_{1s}R_{2s}r^2\,dr \propto \int_0^\infty (2 - r)r^2\mathrm e^{-3r/2}\,dr =
  2 \times \frac{2!}{(3/2)^3} - \frac{3!}{(3/2)^4} = \frac{32}{27} - \frac{32}{27} = 0 .
\]
\end{solution}

\begin{solution}{exo:b2:atomic-orbitals:10}
$\int_0^\infty N^2r^2\mathrm e^{-r}r^2\,dr = N^2 \times 4! = 1$: $N = 1/\sqrt{24} =
1/(2\sqrt6)$, as in the table.
\end{solution}

\begin{solution}{exo:b2:atomic-orbitals:11}
$\langle r\rangle = 1.5a_0$, \omterm{def:b2:atomic-orbitals:radial-distribution}{most probable radius} $a_0$. $P(r) = 4r^2\mathrm e^{-2r}$ rises
steeply from zero and falls slowly: the long tail at large $r$ weighs more in the
mean than in the position of the maximum.
\end{solution}

\begin{solution}{exo:b2:atomic-orbitals:12}
$Y_{p_z} \propto z/r = \cos\theta$ vanishes for $\theta = \pi/2$: the $xy$ plane.
$\psi$ has the sign of $z$, positive above, negative below. On the plane the
\omterm{def:b2:atomic-orbitals:wavefunction}{probability density} is zero.
\end{solution}

\begin{solution}{pb:b2:atomic-orbitals:1}
\textbf{1.} $\int|\psi|^2\,dV = \frac{4\pi}{\pi a_0^3}\int_0^\infty r^2\mathrm e^{-2r/a_0}\,dr =
\frac{4}{a_0^3} \times \frac{a_0^3}{4} = 1$.
\textbf{2.} $P(r) = 4\pi r^2|\psi|^2 = (4/a_0^3)\,r^2\mathrm e^{-2r/a_0}$.
\textbf{3.} $dP/dr = 0$: $2/r - 2/a_0 = 0$, $r = a_0$.
\textbf{4.} $\langle r\rangle = (4/a_0^3)\int_0^\infty r^3\mathrm e^{-2r/a_0}\,dr = (4/a_0^3)
\times 6(a_0/2)^4 = 1.5a_0$.
\textbf{5.} The distribution is skewed, with a long tail at large $r$.
\textbf{6.} $|\psi|^2$ is a density per unit volume, largest at the nucleus; $P(r)$
counts the volume of the shell, $4\pi r^2\,dr$, which vanishes at $r = 0$.
\textbf{7.} \qty{52.9}{pm}.
\textbf{8.} With $x = r/a_0$: $\int_{x_0}^\infty 4x^2\mathrm e^{-2x}\,dx$; by parts twice,
$= \mathrm e^{-2x_0}(2x_0^2 + 2x_0 + 1)$.
\textbf{9.} $5\mathrm e^{-2} = 0.677$.
\textbf{10.} $13\mathrm e^{-4} = 0.238$.
\textbf{11.} $\mathrm e^{-2x}(1 + 2x + 2x^2) = 0.10$ for $x \approx 2.66$: \qty{141}{pm}.
\textbf{12.} The density never vanishes at any distance: no edge. But 90~\% of the
probability lies within \qty{141}{pm}: a practical size.
\textbf{13.} 2s: one \omterm{def:b2:atomic-orbitals:node}{radial node} (sphere of radius $2a_0$), no \omterm{def:b2:atomic-orbitals:node}{angular node}. 2p: one
\omterm{def:b2:atomic-orbitals:node}{angular node} (a plane), no \omterm{def:b2:atomic-orbitals:node}{radial node}.
\textbf{14.} At $r = 2a_0$.
\textbf{15.} $4a_0$ (exercise 2).
\textbf{16.} $d\ln P/dr = 2/r - 2/(2 - r) - 1 = 0$ gives $r^2 - 6r + 4 = 0$, $r = 3 \pm \sqrt5$:
$0.76a_0$ (small inner maximum) and $5.24a_0$ (main maximum).
\textbf{17.} 2s, through its inner maximum. In a many-electron atom a 2s electron
penetrates the 1s shell, is less screened than a 2p electron and lies lower in
energy.
\textbf{18.} All radii are divided by $Z = 6$: 2p at $0.67a_0$ (\qty{35}{pm}); 2s maxima
at $0.13a_0$ and $0.87a_0$.
\textbf{19.} C 3.25, N 3.90, O 4.55.
\textbf{20.} $\varepsilon = -13.6\,Z^{*2}/4$: C \qty{-35.9}{eV}, N \qty{-51.7}{eV}, O
\qty{-70.4}{eV}.
\textbf{21.} As in the chapter: \qty{11.5}{eV} (measured 11.26).
\textbf{22.} N: atom $5 \times (-13.6 \times 3.90^2/4) = \qty{-258.6}{eV}$; \ce{N+}, $Z^* =
4.25$, $4 \times (-13.6 \times 4.25^2/4) = \qty{-245.7}{eV}$; $I = \qty{12.9}{eV}$.
\textbf{23.} O: atom $6 \times (-13.6 \times 4.55^2/4) = \qty{-422.3}{eV}$; \ce{O+}, $Z^* =
4.90$, $5 \times (-13.6 \times 4.90^2/4) = \qty{-408.2}{eV}$; $I = \qty{14.2}{eV}$. Model
11.5, 12.9, 14.2; measured 11.26, 14.53, 13.62.
\textbf{24.} The measured energy of oxygen is lower than that of nitrogen: in
oxygen the fourth 2p electron must share an orbital with another, and their
repulsion makes it easier to remove, while nitrogen's three unpaired p electrons
form a stable half-filled subshell. \omterm{def:b2:atomic-orbitals:slater}{Slater's rules} treat all electrons of a group
alike and know nothing of pairing.
\textbf{25.} The 1s electron of hydrogen is farther than $2a_0$ with probability
$\boldsymbol{13\mathrm e^{-4} \approx 0.238}$.
\end{solution}
